\documentclass[11pt]{article}
\usepackage[T1]{fontenc}
\usepackage[utf8]{inputenc}
\usepackage{amsmath,amssymb}
\usepackage[hidelinks]{hyperref}
\usepackage[margin=1in]{geometry}

\newcommand{\F}{\mathcal F}
\newcommand{\Lip}{\operatorname{Lip}}
\newcommand{\ext}{\operatorname{ext}}
\newcommand{\Mol}{\operatorname{Mol}}

\title{Literature Status Note: Extreme Points of Lipschitz-Free Spaces}
\author{fa\_banach\_001, lane 19}
\date{2026-06-20}

\begin{document}
\maketitle

\paragraph{Status.}
\textbf{literature\_already\_answered}.  This is a literature-status packet, not
an original proof packet.

\paragraph{Source/trigger.}
Ch. Cobollo and A. Peris, \emph{On disjoint dynamical properties and
Lipschitz-free spaces}, arXiv:2405.12626v2.  The extracted signal occurs in
Section 4.1, ``Lipschitz-free working preliminaries'', around source line 635:
the authors recall the area-level conjecture that every extreme point in a
Lipschitz-free space should be a molecule.  This is background context rather
than a numbered new problem of that paper.

\paragraph{Answer source.}
R. J. Aliaga, E. Pernecka, and R. J. Smith, \emph{A solution to the extreme
point problem and other applications of Choquet theory to Lipschitz-free
spaces}, arXiv:2412.04312.  The abstract states that all extreme points of the
unit ball of Lipschitz-free spaces are elementary molecules.  Section 4 is
called ``The extreme point problem''; it records the historical problem and
then gives Theorem 4.1.

\paragraph{Answer theorem.}
Theorem 4.1 of arXiv:2412.04312 states that if \(M\) is a complete metric
space, then
\[
   \ext B_{\F(M)} \subset \Mol(M).
\]
More precisely, the extreme points of \(B_{\F(M)}\) are exactly the molecules
\[
   m_{xy}=\frac{\delta_x-\delta_y}{d(x,y)}
\]
with \(x\neq y\in M\) such that
\[
   d(x,p)+d(p,y)>d(x,y)
   \quad\text{for every }p\in M\setminus\{x,y\}.
\]
Since Lipschitz-free spaces do not change under completion, this supplies the
usual formulation for arbitrary metric spaces after replacing \(M\) by its
completion.

\paragraph{Why this answers the source signal.}
The source signal from arXiv:2405.12626 was exactly the molecule conjecture for
extreme points of Lipschitz-free spaces.  The supporting paper explicitly
announces a solution to that extreme-point problem and provides the final
classification.  Thus this target should not be attacked as a fresh open
problem in this run.

\paragraph{Scope limitation.}
Only the background extreme-point conjecture is cleared.  This packet does not
claim to answer any separate disjoint dynamical property question from
arXiv:2405.12626.

\begin{thebibliography}{9}
\bibitem{CobolloPeris}
Ch. Cobollo and A. Peris,
\emph{On disjoint dynamical properties and Lipschitz-free spaces},
arXiv:2405.12626v2, 2024.

\bibitem{AliagaPerneckaSmith}
R. J. Aliaga, E. Pernecka, and R. J. Smith,
\emph{A solution to the extreme point problem and other applications of Choquet
theory to Lipschitz-free spaces},
arXiv:2412.04312, 2024.
\end{thebibliography}

\end{document}
